\section{Introduction}
\label{sec:intro}

A conversion-rate model must be trained on events that have not happened yet. The click is
observable at once; the conversion it may cause arrives minutes, hours or days later. Wait for
the attribution window to close and the model is always stale by its width. Do not wait, and
every conversion still in flight is recorded as a negative.

A long line of work corrects this, from an explicit model of the delay
distribution~\cite{chapelle2014modeling} through importance weighting~\cite{yasui2020feedback},
elapsed-time sampling~\cite{yang2021capturing}, re-ingesting an event once its outcome is
known~\cite{ktena2019addressing,gu2021real}, unbiased treatment of the fake-negative
population~\cite{chen2022asymptotically}, multiple prediction
horizons~\cite{li2021follow,liu2024online,zhang2025personalized}, and corrective updates from
influence functions~\cite{ding2026delayed}. The statistics differ; but realised as streaming
pipelines, many of them need the training stream to be something one can go back and add to.

We are interested in the case where it is not. In an append-only log a record is written
once, at a fixed age after the click, carrying whatever the label is then, and never rewritten.
Revisiting such a click means either mutating its label --- which turns a stream into a
database and makes a reader's view depend on when it read --- or appending a second record.
Appending is cheaper to build but not free statistically: a click whose conversion arrives late
then appears more often than one that never converts, so the fitted example distribution no
longer matches the traffic, and the multiplier hits training cost at every step. Call the property that avoids both
\emph{single-visit} --- at most one training record per click, its label frozen when written.
It is a restriction one may reasonably pay to escape; we ask only that the price be visible.

We reproduce \nScreened{} delayed-feedback methods from the literature from scratch on the
public Criteo conversion logs and count the training records each one's served model consumes
per click, at a one-day attribution window. \nFail{} write more than one --- up to \maxDup{}
records per click on average, and \maxVisits{} for a single click at worst. Only \passList{}
are single-visit.

What survives is narrow enough to be informative: once a method may write one record, the
only freedom left is \emph{what label it carries} --- a hard label that is wrong (the censored
outcome), one that is right but late (wait for the window), or a target partly or wholly
supplied by a model. We instantiate four points on this \emph{target-hardness} axis and ask
what each degree of softness buys at fixed single-visit cost.

Waiting beats the censored label by a wide margin at short write delays, so the status quo is no
straw man. Replacing part of the hard label with an offline teacher's estimate lowers $\NE$
against waiting in all \cdWaitN{} cells (\cdWaitSig{} significantly); removing the hard label
\emph{entirely} does not, helping only at the smallest capacities and losing to waiting at the
largest. And the hybrid matches an online label corrector~\cite{wang2023unbiased} to within noise ---
which matters less for accuracy than for where each puts the machinery. A model serving several
conversion objectives must reconcile attribution windows differing by orders of magnitude. An
online corrector entangles each window with the served model's own training loop; a teacher
moves it upstream, leaving the served pipeline to consume records that carry a number and no
window at all. We also measure how much of the teacher must stay live: freezing it after initialisation, so
that it becomes a build-time artifact rather than a running service, costs \frzMean{} $\NE$ on
average. This is a precondition for consolidating many objectives, not a demonstration of it;
Sect.~\ref{sec:allinone} states what the evidence supports.
